\section{Results}
\label{sec:results}

\subsection{The corridor where every mechanism is identical}
\label{sec:corridor}

On the replication benchmark, shrinking the catalog from nine classes to six raises aligned macro
F1 from \num{0.4895} to \num{0.7342}, an apparent gain of \num{0.2447}. Over that entire range all
four series --- retirement by F1, retirement by support, confidence rejection, and the control that
applies no mechanism at all --- return the \emph{same} value, and coverage stays at
\num{1.000}.

The reason is arithmetic and it is what makes the observation conclusive. The three classes
retired first hold two, seven and \num{289} parcels out of \num{60000}; the unrestricted argmax
essentially never predicts them, so removing them from the legend removes no parcel from the
delivered map. No mechanism does anything. The metric rises by almost a quarter of a point on its
own.

\begin{table}[tbp]
\caption{Decomposition of the apparent gain on the replication benchmark, nine-class universe.
``Denominator'' is the control of Sect.~\ref{sec:control}: the untouched predictor scored over the
same legend, delivering everything. $K = 5$ is the pre-registered principal criterion.}
\label{tab:decomposition}
\centering
\begin{tabular}{rrrrrr}
\toprule
$K$ & aligned F1 & apparent gain & denominator & mechanism & coverage \\
\midrule
9 & 0.4895 & ---     & ---              & ---            & 1.000 \\
8 & 0.5507 & +0.0612 & +0.0612 (100.0\%) & +0.0000 (0.0\%) & 1.000 \\
7 & 0.6293 & +0.1398 & +0.1398 (100.0\%) & +0.0000 (0.0\%) & 1.000 \\
6 & 0.7342 & +0.2447 & +0.2447 (100.0\%) & +0.0000 (0.0\%) & 1.000 \\
\textbf{5} & \textbf{0.8055} & \textbf{+0.3160} & \textbf{+0.3004 (95.1\%)} & \textbf{+0.0156 (4.9\%)} & \textbf{0.859} \\
4 & 0.8579 & +0.3684 & +0.3352 (91.0\%) & +0.0333 (9.0\%) & 0.806 \\
3 & 0.9126 & +0.4232 & +0.3826 (90.4\%) & +0.0406 (9.6\%) & 0.444 \\
\bottomrule
\end{tabular}
\end{table}

\subsection{The decomposition transports}
\label{sec:transport}

The same protocol, unchanged, gives the same qualitative answer on the primary benchmark. Shrinking
from eighteen classes to eight raises aligned macro F1 from \num{0.5612} to \num{0.8052}, an
apparent gain of \num{0.2440}, of which the control recovers \num{0.2155}, or
\num{88.3}\,\%. At the pre-registered operating point of $K = 9$ the control recovers
\num{87.1}\,\% and the result holds for both of the two predictors we tested.

Two benchmarks, two years, two regions, two model families and one protocol therefore agree: the
denominator dominates. Where they differ is in degree, and in the direction that matters for a
practitioner: the replication benchmark, whose imbalance is an order of magnitude worse, leaves the
mechanism even less to explain.

\begin{figure}[tbp]
\centering
\includegraphics[width=\textwidth]{frontier-primary.pdf}
\caption{The same frontier on the primary benchmark, for its two strongest predictors. The control
that applies no mechanism (crosses) already recovers most of the distance travelled by the
mechanisms, and retirement by support (squares) again trails the other two throughout.}
\label{fig:primary}
\end{figure}

\begin{figure}[tbp]
\centering
\includegraphics[width=\textwidth]{frontier-replication.pdf}
\caption{Quality--coverage frontier on the replication benchmark, in the two pre-registered class
universes. The four series are indistinguishable until $K = 6$ because coverage is complete and no
mechanism acts; the dotted line marks the principal criterion. Retirement by support (squares) is
the only series that falls as more classes are retired.}
\label{fig:replication}
\end{figure}

\subsection{At equal coverage, abstention wins}
\label{sec:ordering}

Once coverage falls below one and the mechanisms begin to differ, they can be ordered. At the
principal criterion of the replication benchmark, confidence rejection reaches \num{0.8303} against
\num{0.8055} for legend shrinking at the same coverage, a difference of \num{-0.0248} with a paired
\num{95}\,\% interval of $(-0.0272, -0.0225)$ and a Holm-adjusted $p$ below \num{0.0001}. Both
regions agree in sign and magnitude, at \num{-0.0248} and \num{-0.0249}.

On the primary benchmark the difference has the same sign but is not distinguishable from zero:
\num{-0.0288} with an interval of $(-0.0414, +0.0080)$ and a Holm-adjusted $p$ of \num{0.336}. The
honest summary across both benchmarks is therefore that we find no evidence that shrinking the
legend outperforms abstention, and on the larger and more imbalanced benchmark we find significant
evidence of the reverse.

\begin{figure}[tbp]
\centering
\includegraphics[width=\textwidth]{cobertura.pdf}
\caption{The same measurements plotted against delivered coverage rather than legend size, which
is how the selective classification literature reads them. The control (crosses) sits on the
coverage-one vertical by construction, since it never withholds anything; the point of the figure
is how many of those crosses lie \emph{above} the mechanisms' curves at lower coverage.}
\label{fig:coverage}
\end{figure}

\subsection{Retirement by support is the weakest criterion, and it is not monotone}
\label{sec:support}

Retiring classes by low support is the criterion practitioners describe using, and it is the worst
of the three on both benchmarks. On the replication benchmark it reaches \num{0.7129} at the
principal criterion against \num{0.8055} for retirement by F1, a gap of \num{0.0926}, which widens
to \num{0.2085} at $K = 3$. On the primary benchmark the corresponding gap is between \num{0.12}
and \num{0.13}.

It also does something neither other criterion does: it \emph{degrades} as more classes are
retired, from \num{0.7129} at $K = 5$ to \num{0.7041} at $K = 3$. Inspecting the legends explains
why. Support-based retirement removes rapeseed, which holds \num{1718} parcels and which the model
separates well, while retaining both meadow classes, which the model confuses with each other. The
criterion optimizes sample size, which is not the same quantity as the quality of what is
delivered. Fig.~\ref{fig:legends} shows the substitution directly.

\begin{figure}[tbp]
\centering
\includegraphics[width=\textwidth]{leyendas.pdf}
\caption{Which classes each criterion keeps as the legend shrinks, on the replication benchmark;
a filled cell means the class is promised. Support-based retirement (right) drops rapeseed at
$K = 5$ while keeping both meadow classes, which the model confuses with one another;
quality-based retirement (left) does the reverse. The two criteria disagree about exactly the
classes that decide the outcome.}
\label{fig:legends}
\end{figure}

\subsection{A trained meta-learner retires classes implicitly}
\label{sec:arbiter}

The mechanism this paper formalizes is already being applied, undeclared, by trained ensembles. On
the primary benchmark, a stacking meta-learner over three members drives one class of \num{355}
parcels to an F1 of exactly \num{0.0000}, while the best single model reaches \num{0.7790} on that
same class. The meta-learner has retired a class from the product without anyone choosing to, and
without it appearing anywhere in a results table.

\begin{table}[tbp]
\caption{Per-class F1 on the primary benchmark, leak-free, ordered by support. The trained
arbiter is the only combiner that zeroes a class, and it does so to one of \num{355} parcels that
every other combiner --- including the untrained mean --- handles above \num{0.77}. The retirement
is therefore a learned behavior, not a property of combining.}
\label{tab:perclass}
\centering
\small
\begin{tabular}{lrrrrr}
\toprule
Class & support & best single & mean & weighted vote & trained arbiter \\
\midrule
Meadow & 6\,128 & 0.9099 & 0.8980 & 0.9046 & 0.9138 \\
Corn & 2\,196 & 0.9295 & 0.8683 & 0.9053 & 0.9274 \\
Grapevine & 2\,067 & 0.9066 & 0.9141 & 0.9154 & 0.9124 \\
Soft winter wheat & 1\,584 & 0.8922 & 0.8466 & 0.8712 & 0.8886 \\
Leguminous fodder & 660 & 0.5087 & 0.3498 & 0.4707 & 0.5243 \\
Winter rapeseed & 387 & 0.9358 & 0.9153 & 0.9275 & 0.9377 \\
\textbf{Winter durum wheat} & \textbf{355} & \textbf{0.7790} & \textbf{0.7752} & \textbf{0.7782} & \textbf{0.0000} \\
Spring barley & 198 & 0.6667 & 0.4691 & 0.5879 & 0.5228 \\
Mixed cereal & 193 & 0.4949 & 0.3525 & 0.4302 & 0.2709 \\
Potatoes & 103 & 0.5595 & 0.5143 & 0.5638 & 0.5823 \\
\bottomrule
\end{tabular}
\end{table}

\begin{table}[tbp]
\caption{The same stacking meta-learner under three estimators. These are not three plausible
choices: the first is contaminated and the third is penalized by an aggregation artifact.}
\label{tab:regimes}
\centering
\begin{tabular}{lrl}
\toprule
Estimator & macro F1 & what it measures \\
\midrule
Refit on all parcels & 0.7470 & in-sample for the meta-learner \\
Pooled spatial out-of-fold & 0.6789 & leak-free, scored once \\
Mean of per-block macros & 0.5360 & leak-free, penalized by aggregation \\
\bottomrule
\end{tabular}
\end{table}

The regime in which such a meta-learner is measured is decisive. Scored in-sample --- refitted on
all parcels after cross-validation, which is the default of many stacking implementations --- it
reports \num{0.7470}. Scored leak-free by pooling its spatial out-of-fold posteriors and scoring
once, it reports \num{0.6789}. Averaging per-block macro scores instead of pooling reports
\num{0.5360}. The leakage accounts for \num{0.068} and the choice of aggregation for a further
\num{0.143}; conflating the two leads to reporting \num{0.536} as the honest figure when it is
\num{0.679}, penalizing the model for a measurement artifact on top of its actual error. Against
the best single model, the leak-free meta-learner loses by \num{0.0572} with an interval of
$(-0.0644, -0.0499)$ that excludes zero.
